\editionPage[\pdfbookmark[0]{Corrigés des automatismes}{bm-automatismes-cor-1-parent}]{automatismes-cor-1}{Automatismes · Corrigés}{Séries A et B · Nombres, calculs et proportions}{\hyperref[automatismes-1]{Énoncés p.\,\pageref*{automatismes-1}} · Réponses et justifications courtes}
\begin{blocEdition}{automatismes-cor-1}
\editionEtape{Série A · Nombres, calculs et grandeurs}
\begin{deuxcolonnes}\begin{questions}
\item \begin{minipage}[t]{\linewidth}\raggedright\leavevmode \otEtiq{otMuted}{Calcul}\enspace $\dfrac8{12}+\dfrac9{12}=\res{\dfrac{17}{12}}$ ; $\dfrac{25}{40}>\dfrac{24}{40}$, donc $\res{\dfrac58>\dfrac35}$.\end{minipage}
\item \begin{minipage}[t]{\linewidth}\raggedright\leavevmode \otEtiq{otMuted}{Calcul}\enspace $2^{5-2-4}=\res{2^{-1}}=0{,}5$.\end{minipage}
\item \begin{minipage}[t]{\linewidth}\raggedright\leavevmode \otEtiq{otMuted}{Calcul}\enspace $0{,}35=\dfrac{35}{100}=\res{\dfrac7{20}}=\res{35\,\%}$.\end{minipage}
\item \begin{minipage}[t]{\linewidth}\raggedright\leavevmode \otEtiq{otMuted}{Calcul}\enspace $50\times0{,}2=\res{10}$ ; le produit exact vaut 10,437.\end{minipage}
\item \begin{minipage}[t]{\linewidth}\raggedright\leavevmode \otEtiq{otMuted}{Calcul}\enspace $\dfrac ab=\dfrac{1001\times999}{1000^2}=\dfrac{999\,999}{1\,000\,000}<1$ et $b>0$, donc $\res{a<b}$.\end{minipage}
\item \begin{minipage}[t]{\linewidth}\raggedright\leavevmode \otEtiq{otMuted}{Calcul}\enspace $90\,000/3600=\res{25\text{ m/s}}$ ; $2{,}5\times10\,000=\res{25\,000\text{ cm}^2}$.\end{minipage}
\item \begin{minipage}[t]{\linewidth}\raggedright\leavevmode \otEtiq{otMuted}{Calcul}\enspace $-(3-x)=\res{x-3}$ ; diviser par $\dfrac34$ revient à multiplier par $\dfrac43$ : $\res{\dfrac83}$.\end{minipage}
\item \begin{minipage}[t]{\linewidth}\raggedright\leavevmode \otEtiq{otMuted}{Évolutions}\enspace $1-0{,}08=\res{0{,}92}$ ; multiplier par 1,15, c’est $\res{\text{augmenter de }15\,\%}$.\end{minipage}
\item \begin{minipage}[t]{\linewidth}\raggedright\leavevmode \otEtiq{otMuted}{Probabilités}\enspace C’est $\res{\text{impossible}}$ : une probabilité est comprise entre 0 et 1. Il faut reprendre le calcul.\end{minipage}
\item \begin{minipage}[t]{\linewidth}\raggedright\leavevmode \otEtiq{otMuted}{Géométrie}\enspace $\cos\widehat{ABC}=\dfrac{AB}{BC}=\res{\dfrac12}$, donc $\res{\widehat{ABC}=60^\circ}$.\end{minipage}
\end{questions}\end{deuxcolonnes}
\editionEtape{Série B · Calcul littéral, équations et proportions}
\begin{deuxcolonnes}\begin{questions}
\item \begin{minipage}[t]{\linewidth}\raggedright\leavevmode \otEtiq{otMuted}{Calcul}\enspace $\res{4x^2+12x+9}$ ; le double produit vaut $12x$.\end{minipage}
\item \begin{minipage}[t]{\linewidth}\raggedright\leavevmode \otEtiq{otMuted}{Calcul}\enspace $\res{5x(x-3)}$ ; $\res{(x-7)(x+7)}$.\end{minipage}
\item \begin{minipage}[t]{\linewidth}\raggedright\leavevmode \otEtiq{otMuted}{Calcul}\enspace $\res{S=\{-6;6\}}$ : un nombre et son opposé ont le même carré.\end{minipage}
\item \begin{minipage}[t]{\linewidth}\raggedright\leavevmode \otEtiq{otMuted}{Calcul}\enspace $2x=10$, donc $\res{x=5}$.\end{minipage}
\item \begin{minipage}[t]{\linewidth}\raggedright\leavevmode \otEtiq{otMuted}{Calcul}\enspace $6=2x$, donc $\res{x=3}$, valeur autorisée.\end{minipage}
\item \begin{minipage}[t]{\linewidth}\raggedright\leavevmode \otEtiq{otMuted}{Calcul}\enspace $-2x>-4$ ; diviser par $-2$ renverse le sens : $\res{x<2}$.\end{minipage}
\item \begin{minipage}[t]{\linewidth}\raggedright\leavevmode \otEtiq{otMuted}{Calcul}\enspace $\res{R=\dfrac UI}$ ; $\res{h=\dfrac{3V}B}$.\end{minipage}
\item \begin{minipage}[t]{\linewidth}\raggedright\leavevmode \otEtiq{otMuted}{Calcul}\enspace $\dfrac12\times2\times3^2=\res{9\text{ J}}$.\end{minipage}
\item \begin{minipage}[t]{\linewidth}\raggedright\leavevmode \otEtiq{otMuted}{Proportions}\enspace $9/0{,}30=\res{30\text{ élèves}}$.\end{minipage}
\item \begin{minipage}[t]{\linewidth}\raggedright\leavevmode \otEtiq{otMuted}{Proportions}\enspace $0{,}15\times240=\res{36}$ ; $15\,\%=0{,}15=\res{\dfrac3{20}}$.\end{minipage}
\end{questions}\end{deuxcolonnes}

\end{blocEdition}
\editionPage{automatismes-cor-2}{Automatismes · Corrigés}{Séries C et D · Géométrie, fonctions, données et hasard}{\hyperref[automatismes-2]{Énoncés p.\,\pageref*{automatismes-2}} · Réponses et justifications courtes}
\begin{blocEdition}{automatismes-cor-2}
\editionEtape{Série C · Géométrie et mesures}
\begin{deuxcolonnes}\begin{questions}
\item \begin{minipage}[t]{\linewidth}\raggedright\leavevmode \otEtiq{otMuted}{Géométrie}\enspace Chaque unité est partagée en 4 : un pas vaut $0{,}25$. $\res{A(-1{,}25)}$ et $\res{B(0{,}75)}$.\end{minipage}
\item \begin{minipage}[t]{\linewidth}\raggedright\leavevmode \otEtiq{otMuted}{Géométrie}\enspace $\res{6\pi\text{ cm}}$ et $\res{9\pi\text{ cm}^2}$.\end{minipage}
\item \begin{minipage}[t]{\linewidth}\raggedright\leavevmode \otEtiq{otMuted}{Géométrie}\enspace $7\times4/2=\res{14\text{ cm}^2}$.\end{minipage}
\item \begin{minipage}[t]{\linewidth}\raggedright\leavevmode \otEtiq{otMuted}{Géométrie}\enspace $\res{30\text{ cm}^3}$, puis $30/3=\res{10\text{ cm}^3}$.\end{minipage}
\item \begin{minipage}[t]{\linewidth}\raggedright\leavevmode \otEtiq{otMuted}{Géométrie}\enspace $\pi\times2^2\times5=\res{20\pi\text{ cm}^3}$ ; $\res{\dfrac{20\pi}3\text{ cm}^3}$.\end{minipage}
\item \begin{minipage}[t]{\linewidth}\raggedright\leavevmode \otEtiq{otMuted}{Géométrie}\enspace $\res{60\text{ cm}^3}$ ; $\dfrac43\pi\times3^3=\res{36\pi\text{ cm}^3}$.\end{minipage}
\item \begin{minipage}[t]{\linewidth}\raggedright\leavevmode \otEtiq{otMuted}{Géométrie}\enspace $\sqrt{36+64}=\res{10\text{ cm}}$ (Pythagore).\end{minipage}
\item \begin{minipage}[t]{\linewidth}\raggedright\leavevmode \otEtiq{otMuted}{Géométrie}\enspace Thalès : $\dfrac{MN}{BC}=\dfrac{AM}{AB}=\dfrac13$, donc $\res{MN=3}$.\end{minipage}
\item \begin{minipage}[t]{\linewidth}\raggedright\leavevmode \otEtiq{otMuted}{Géométrie}\enspace $10\sin30^\circ=10\times\dfrac12=\res5$.\end{minipage}
\item \begin{minipage}[t]{\linewidth}\raggedright\leavevmode \otEtiq{otMuted}{Calcul}\enspace $\res{1500\text{ cm}^3}$ ; $2+\dfrac{15}{60}=\res{2{,}25\text{ h}}$.\end{minipage}
\end{questions}\end{deuxcolonnes}
\editionEtape{Série D · Fonctions, statistiques et probabilités}
\begin{deuxcolonnes}\begin{questions}
\item \begin{minipage}[t]{\linewidth}\raggedright\leavevmode \otEtiq{otMuted}{Fonctions}\enspace Le point de $\mathcal C_f$ d’abscisse 0 est $(0;-1)$ : $\res{f(0)=-1}$.\end{minipage}
\item \begin{minipage}[t]{\linewidth}\raggedright\leavevmode \otEtiq{otMuted}{Fonctions}\enspace La droite $y=1$ coupe $\mathcal C_f$ en $x=-1$ et $x=2$ : $\res{-1\text{ et }2}$.\end{minipage}
\item \begin{minipage}[t]{\linewidth}\raggedright\leavevmode \otEtiq{otMuted}{Géométrie}\enspace $\res{M(3;-2)}$ : abscisse d’abord, ordonnée ensuite.\end{minipage}
\item \begin{minipage}[t]{\linewidth}\raggedright\leavevmode \otEtiq{otMuted}{Fonctions}\enspace $2^2+3=7$ : $\res{\text{oui}}$. Pour $x=-1$ : $\res{y=4}$.\end{minipage}
\item \begin{minipage}[t]{\linewidth}\raggedright\leavevmode \otEtiq{otMuted}{Fonctions}\enspace $f$ est $\res{\text{linéaire}}$ (donc affine), $g$ est $\res{\text{affine}}$ ; leurs courbes sont des droites. $h$ n’est pas affine.\end{minipage}
\item \begin{minipage}[t]{\linewidth}\raggedright\leavevmode \otEtiq{otMuted}{Statistiques}\enspace Le demi-disque vaut $180^\circ$ : $\dfrac{45}{180}=\res{\dfrac14=25\,\%}$.\end{minipage}
\item \begin{minipage}[t]{\linewidth}\raggedright\leavevmode \otEtiq{otMuted}{Statistiques}\enspace Moyenne $24/4=\res6$ ; médiane $(5+6)/2=\res{5{,}5}$.\end{minipage}
\item \begin{minipage}[t]{\linewidth}\raggedright\leavevmode \otEtiq{otMuted}{Statistiques}\enspace $\dfrac{3\times5+7\times15}{10}=\res{12}$.\end{minipage}
\item \begin{minipage}[t]{\linewidth}\raggedright\leavevmode \otEtiq{otMuted}{Statistiques}\enspace Même médiane 9. Écarts interquartiles 6 et 3 : $\res{B\text{ est moins dispersée}}$ autour de sa médiane.\end{minipage}
\item \begin{minipage}[t]{\linewidth}\raggedright\leavevmode \otEtiq{otMuted}{Probabilités}\enspace Issues 3, 6, 9, chacune de probabilité $\dfrac1{10}$ : $P(A)=\dfrac{\operatorname{Card}(A)}{\operatorname{Card}(\Omega)}=\res{\dfrac3{10}}$ ; contraire : $1-\dfrac3{10}=\res{\dfrac7{10}}$.\end{minipage}
\end{questions}\end{deuxcolonnes}

\end{blocEdition}
